\chapter{Ziel und Einordnung}

\section{Anschluss statt Wiederholung}

Band~3 liefert die Mengenaxiome; die Bände über Funktionen, Bijektionen,
Wohlordnungen und reelle Zahlen liefern die bisher entwickelten Begriffe.
Dieser Band setzt diese Ergebnisse voraus. Neu ist die Untersuchung der
\emph{Reichweite} der Axiome: Welche Aussagen folgen aus ihnen, und welche
werden durch sie nicht entschieden? Die Kapitel~4 und~5 bleiben die
Zielkapitel. Die dafür erforderlichen Konstruktionen werden in den
Kapiteln~2 und~3 aufgebaut, nicht auf einen weiteren Band verschoben.

\FormulaDefDeltaK[Grundtheorie dieses Bandes]
{\BXLVIIIZFC\coloneqq\BXLVIIIZF+\mathsf{AC}}
{B48BaseTheory}{
  \DeltaRow{Theorie}{\BXLVIIIZF}
  \DeltaRow{Zusatz}{\mathsf{AC}}
}
\noindent
Mit \(\BXLVIIIZF\) sind Extensionalität, leere Menge, Paarmenge,
Vereinigung, Potenzmenge, Unendlichkeit, Fundierung sowie die
Aussonderungs- und Ersetzungsschemata gemeint; \(\mathsf{AC}\) ist das
Auswahlaxiom. Dies sind die bereits eingeführten Grundlagen, keine zweite
Liste neu zu rechtfertigender Mengenaxiome. Insbesondere wird das
Ersetzungsschema \FormulaRefAuto{ReplacementSchema}[ax] verwendet.
Ein Schema hat für jede Formel eine eigene Instanz.

Der axiomatische Aufbau besteht hier aus Definitionen, danach bewiesenen
Hilfssätzen und schließlich Anwendungen. Weder die Existenz des
konstruktiblen Universums noch der Forcingerweiterungen wird als neues
Axiom eingeführt. Ein Modell mit GCH und ein Modell mit einem geeigneten
Plateau der Potenzfunktion werden aus einer Konsistenzannahme gewonnen.

\section{Die beiden Zielaussagen}

Wir schreiben \(A\approx B\), wenn eine Bijektion zwischen \(A\) und \(B\)
existiert, und \(\mathcal P(A)\) für die Potenzmenge von \(A\).

\FormulaDefDeltaK[Globale Potenzmengenkürzung]
{\BXLVIIIPC\;\coloneqq\;
 \forall A\,\forall B\,
 \bigl(\mathcal P(A)\approx\mathcal P(B)\ \longrightarrow\ A\approx B\bigr)}
{B48PCStatement}{
  \DeltaRow{Mengen}{A,B}
}

\noindent
Die Quantoren sind wesentlich: Untersucht wird eine einzige Aussage über
\emph{alle} Mengen, nicht die Entscheidbarkeit eines beliebigen konkreten
Paares \(A,B\). Für endliche Mengen gilt die Kürzung. Für beliebige Mengen
ist sie unter der Konsistenzvoraussetzung unabhängig von ZFC.

Davon zu unterscheiden ist ein Ergebnis aus Band~8:
Ist die durch eine bestimmte Abbildung \(f\) induzierte Potenzmengenabbildung
bijektiv, so ist bereits \(f\) bijektiv; siehe
\FormulaRefAuto{PowerMapBijectiveReflectsBijectivity}[thm].
Eine \emph{beliebige} Bijektion zwischen zwei Potenzmengen muss aber nicht
von einer Abbildung zwischen den Ausgangsmengen induziert sein.

Die zweite Zielaussage ist
\[
 \BXLVIIICH:\quad 2^{\aleph_0}=\aleph_1.
\]
Die Identifikation \(2^{\aleph_0}=|\mathbb R|\) und die Äquivalenz mit der
Aussage, es gebe keine dazwischenliegende Mächtigkeit, werden in
Kapitel~5 bewiesen.

\section{Objektsprache und Metasprache}

\FormulaDefDeltaK[Konsistenz und Unabhängigkeit]
{\begin{aligned}
 \BXLVIIICon(T)&\coloneqq \neg(T\vdash\bot),\\
 \operatorname{Unabh}(T,\varphi)&\coloneqq
 (T\nvdash\varphi)\land(T\nvdash\neg\varphi)
\end{aligned}}
{B48ConsistencyNotation}{
  \DeltaRow{Theorie}{T}
  \DeltaRow{Satz}{\varphi}
}

\noindent
\(T\vdash\varphi\) besagt, dass eine endliche formale Herleitung aus \(T\)
existiert. \(M\models\varphi\) besagt dagegen, dass \(\varphi\) in einer
Struktur wahr ist. In allen abschließenden Unabhängigkeitssätzen steht
\(\BXLVIIICon(\BXLVIIIZFC)\) ausdrücklich als Voraussetzung. Der Band
behauptet nicht, die Widerspruchsfreiheit von ZFC voraussetzungslos oder
innerhalb von ZFC zu beweisen.

Die Metatheorie ist die gewöhnliche klassische Mengenmathematik auf
ZFC-Basis. In ihr werden Syntax, endliche Beweise und Mengenstrukturen
untersucht. Die Henkin-Konstruktion beschreibt den vorgesehenen Beweis
des Vollstaendigkeitssatzes; die lueckenlose Ausarbeitung ist noch offen.
Der Goedel'sche \emph{Unvollstaendigkeitssatz} ist dafuer keine
Voraussetzung.

\section{Lesart der Beweistabellen}

Die Tabellen verwenden wie die anderen Bände offene Voraussetzungen,
nummerierte Schritte und Begründungen. Ein Verweis auf einen bewiesenen
Satz darf eine bereits geleistete Ableitung zusammenfassen.
Strukturelle und transfinite Induktionen werden mit ihrer
Induktionsbehauptung und den entscheidenden Fällen angegeben.
Es handelt sich um mathematische Beweise in Tabellenform, nicht um eine
maschinell geprüfte Formalisierung jedes logischen Einzelschrittes.

Bei Klassen wie \(L\) ist \(x\in L\) eine Abkürzung für eine Formel, keine
zusätzliche Art von Objekt. Die Formulierung \(L\models\BXLVIIIZFC\)
bedeutet ein \emph{Schema}: Für jedes einzelne gewöhnliche ZFC-Axiom wird
seine Relativierung auf \(L\) bewiesen. Ebenso ist die Forcingrelation
jeweils für eine feste Formel definiert. Eine Wahrheitsdefinition für das
gesamte Mengenuniversum wird an keiner Stelle vorausgesetzt.
